\documentclass[11pt]{article}
\usepackage{amsmath,amssymb,amsthm,geometry,booktabs}
\geometry{margin=1in}
\setlength{\emergencystretch}{1em}
\newtheorem{theorem}{Theorem}
\newcommand{\Aut}{\operatorname{Aut}}
\newcommand{\im}{\operatorname{im}}
\title{The augmentation module of the Steiner system on 23 points}
\author{Working manuscript: independent review pending}
\date{}
\begin{document}\maketitle

This note completes the length-$253$ instance of the common incidence
argument. Its source is Leemans and Rodrigues, \emph{Linear codes with
complementary duals from some strongly regular subgraphs of the McLaughlin
graph}, Math. Commun.\ 21 (2016), 239--249 (L050), Theorem~1(d) and
Proposition~5. We retain the source's Steiner construction and the exact
weight enumerations already completed for this construction. The remaining
irreducibility and full automorphism assertions follow from the same two
certificate arguments used for the length-$120$ and length-$176$ cases.
The finite programs and their literal inputs are necessary proof premises.
No historical novelty or full Lean formalization is asserted.

Let $\mathcal D$ be the specified $S(4,7,23)$, with point set $\Omega$
and $253$ heptads. Its point--block incidence matrix is $H$.
Two distinct heptads intersect in one or three points, and $\Gamma$ joins
them in the latter case. Write $A$ for its adjacency matrix. All vector
spaces in this note are over $F=\mathbb F_3$. Put
\[
 E=\{x\in F^{23}:\textstyle\sum_i x_i=0\},\quad
 C=\im A,\quad D=H^{\mathsf T}E.
\]
The shared construction checks the Steiner condition and actual incidence
matrix, not merely the strongly regular parameters.

\begin{theorem}
The code $C$ is the orthogonal direct sum $\langle\mathbf j\rangle\oplus D$,
where $D$ is an absolutely irreducible $22$-dimensional module for $M_{23}$.
Both $C$ and $D$ are LCD. Their parameters and minimum multiplicities are
\[
\begin{array}{c|c|r}
\text{code}&\text{parameters}&\text{minimum words}\\\hline
C &[253,23,77]_3&46\\
D &[253,22,97]_3&3542\\
C^\perp &[253,230,6]_3&850080\\
D^\perp &[253,231,6]_3&850080.
\end{array}
\]
The full groups of coordinate permutations preserving $C$ and $\Gamma$
are equal and isomorphic to $M_{23}$.
\end{theorem}

\begin{proof}
There are $77$ heptads through each point and $21$ through each pair.
Consequently
\[
 HH^{\mathsf T}=56I_{23}+21J_{23}=2I_{23}\quad\text{over }F.
\]
In particular $H^{\mathsf T}$ is injective. Every column sum is seven,
so $H^{\mathsf T}\mathbf1=\mathbf j$, and the two intersection sizes give
\[
 A=J_{253}-H^{\mathsf T}H
   =H^{\mathsf T}(J_{23}-I_{23})H.
\]
The middle matrix is invertible: its eigenvalues on the constant line
and its sum-zero complement are $22=1$ and $-1$ in $F$.
Thus $C=\im H^{\mathsf T}$, of dimension $23$.
Since $23\ne0$ in $F$, the point space decomposes orthogonally into its
constant line and $E$. The transferred form is twice the ordinary dot
product. It is nondegenerate on both summands. Their images give the
asserted orthogonal splitting and both LCD conclusions.

We next prove the irreducibility, using explicit matrices in place of a
modular character-table inference. The five permutations of $\Omega$
in \texttt{generators.json} preserve the full heptad set. They therefore
act on the code coordinates and on $E$. In the basis
$e_i-e_{23}$, $1\leq i\leq22$, the row matrix for a point permutation
$g$ has entries
\[
 M(g)_{ij}=\mathbf1_{j=g(i)}-\mathbf1_{j=g(23)},\qquad 1\leq i,j\leq22.
\]
An indicator with point $23$ is zero because $j\leq22$.
The certificate checks that multiplication by these matrices agrees
with the induced permutation of all $253$ coordinates on the actual
rows $h_i-h_{23}$. It produces $484$ independent words in these matrices,
each of length at most six, and reevaluates every word before reducing
the flattened matrices. They span $M_{22}(F)$.

Any invariant subspace of $E$ is therefore stable under every endomorphism
of $E$. A nonzero vector can be sent to any prescribed vector by an
endomorphism, so such a subspace must be all of $E$. The same independent
matrix words remain independent after any field extension. This proves
absolute irreducibility of $E$, and hence of its equivariant image $D$.
The argument is the general full-matrix-algebra criterion already used
for the other residual codes.

The existing enumeration records contain the complete weights of the
row-reduced adjacency code and of the difference code. Their bases are
compared byte for byte with the bases reconstructed here. The latter
row space is $D$: its dimension is $22$, it lies in $C$ and is orthogonal
to $\mathbf j$. The recorded enumerations have respectively $3^{23}$
and $3^{22}$ coefficient vectors. They give the two distances and counts
in the first two rows of the table. The dual distributions follow from
the classical ternary MacWilliams identity
\[
 A_j^\perp=3^{-k}\sum_{i=0}^{253} A_i
 \sum_s(-1)^s2^{j-s}\binom is\binom{253-i}{j-s}.
\]
All coefficients are evaluated over the integers. Nonnegativity,
integrality and total $3^{253-k}$ are checked. Both duals have no words
of weights one through five and exactly $850080$ of weight six.
This proves the table without repeating either exhaustive enumeration.

It remains to prove that the displayed automorphism subgroup is full.
Each point-incidence row $h_i$ is a word of $C$ of weight $77$.
The $46$ words $h_i,-h_i$ are distinct. By the complete minimum count
they are all minimum words. Their $23$ supports are exactly the sets of
heptads through the respective points. A code automorphism permutes these
supports and thus induces an incidence-preserving permutation of the
points and blocks of $\mathcal D$. Conversely every design automorphism
permutes the incidence rows and preserves $C$. Hence $\Aut(C)$ is the
group induced on the blocks by $\Aut(\mathcal D)$.

For an explicit upper bound, form the bipartite point--block incidence
graph, with its two parts initially given different colours. Refine the
colours by the old colour and the counts of neighbours of each colour.
After individualizing points $1,2,3,4,5,6$ successively, the certificate
obtains current cell sizes
\[
 23,\quad22,\quad21,\quad20,\quad16,\quad3.
\]
All points are then distinguished. A pointwise automorphism fixes every
heptad, since distinct heptads have different point neighbourhoods.
Each cell contains the orbit of its chosen point under the preceding
stabilizer; repeated orbit--stabilizer therefore gives
\[
 |\Aut(\mathcal D)|\leq23\cdot22\cdot21\cdot20\cdot16\cdot3=10200960.
\]
Exact group computations on the literal generators give a subgroup of
this order and an isomorphism to the standard group $M_{23}$.
The induced block action is checked to be faithful. Thus the upper bound
is attained and identifies the full design and code groups. Finally the
subgroup preserves $\Gamma$, and every graph automorphism preserves its
adjacency row space. Therefore
$G\leq\Aut(\Gamma)\leq\Aut(C)=G$, proving the asserted equality.
The upper bound uses only incidence refinement; graph-isomorphism
software was used to discover the generators and is not a final premise.
\end{proof}

\paragraph{Qualifications to Proposition 5.}
The main Theorem~1(d) is recovered as stated. Proposition~5 additionally
assigns distance eight to $D^\perp$ and declares it irreducible. Both
assertions are false for these codes. Since $D\subset C$, the invariant
space $C^\perp$ of dimension $230$ is a proper nonzero subspace of
$D^\perp$, whose dimension is $231$. The same inclusion already gives
$d(D^\perp)\leq d(C^\perp)=6$; the exact value six is verified above.
The proposition also calls $D$ the ``unique and smallest'' irreducible
module without qualifications. The unqualified statement overlooks the
trivial one-dimensional module. Its intended classification requires an
input separate from our matrix certificate. We retain the complete
$3$-modular character table of $M_{23}$, as in the source's reference [13],
Jansen--Lux--Parker--Wilson, \emph{An Atlas of Brauer Characters} (1995).
Its irreducible degrees over an algebraic closure are
\[
 1,22,45,45,104,104,231,253,770,770,990,990,1035.
\]
The accompanying GAP calculation checks these entries against its
installed table \texttt{M23mod3}; it does not reprove completeness of
that table. Together with standard scalar extension and descent for
simple modules over finite fields, the table gives a unique
$22$-dimensional irreducible $F$-module, and $22$ is the least
nontrivial degree. Indeed a simple module splits after scalar extension
into an orbit of absolutely simple constituents of equal degree. Below
degree $22$ the only available constituent is trivial, which forces the
original simple module to be trivial. At degree $22$ the unique
constituent is the one already realized by $D$; descent of module
isomorphisms gives uniqueness over $F$. These are retained
classification premises, explicitly distinct from the direct proof of
irreducibility of $D$. The source's open uniqueness question for the
strongly regular parameters $(253,140,87,65)$ is not resolved here.

\end{document}
